\begin{table}[tb]
\centering\scriptsize\setlength{\tabcolsep}{2pt}
\caption{全24条件の論理tokenと成否．成：成功，否：未達．$+?_{n}$は別に使用量不明が$n$生成あることを示す．M：残時間表示，H：非表示．}
\begin{tabular}{llrrr}
課題 & 条件 & 反復1（P80） & 反復2 & 反復3 \\ \hline
写真 & M & 177,373 成 & 210,357 成 & 164,024 成 \\
写真 & H & 135,628 成 & 135,628 成 & 135,628 成 \\
精算 & M & 284,294 成 & 215,182 成 & 263,695 成 \\
精算 & H & 211,110 成 & 211,110 成 & 211,110 成 \\
返金 & M & 216,854 成 & 138,701 成 & 73,670 成 \\
返金 & H & 102,087 成 & 102,087 成 & 102,087 成 \\
書出し & M & 108,571 否 & 134,352 否 & 439,709 否 \\
書出し & H & 331,569$+?_{1}$ 否 & 319,348 否 & 324,014$+?_{1}$ 否 \\
\hline
\end{tabular}
\end{table}
反復1はMが4評価中3成功，Hが4評価中3成功で，全論理tokenは順に787,092，780,394$+?_{1}$だった．
全費用差は使用量不明または未実行のため未確定である．
開発上の採用基準は未達で，作業時間合計はM 1,897.7秒，H 2,975.3秒だった．
反復2はMが4評価中3成功，Hが4評価中3成功で，全論理tokenは順に698,592，768,173だった．
H−Mの合計差は+69,581 tokenである．
開発上の採用基準は未達で，作業時間合計はM 1,991.5秒，H 2,691.4秒だった．
反復3はMが4評価中3成功，Hが4評価中3成功で，全論理tokenは順に941,098，772,839$+?_{1}$だった．
全費用差は使用量不明または未実行のため未確定である．
開発上の採用基準は未達で，作業時間合計はM 2,645.3秒，H 2,975.6秒だった．
元のP80にはHの書出し中断による使用量不明１生成がある．事後の補助解析では，単一slotで直列に実行した全143送信を推論器ログへ順序対応させ，既知142生成の入力処理量と出力量を照合した．元の終端使用量を復元したものではないが，同じ送信の事前入力計数と生成開始後の単調な出力計数を使えるという前提では，Hの論理費用は820,051以上となり，Mの787,092を少なくとも32,959上回る．順序対応は要求IDの暗号学的な結合ではなく，未知費用は元表に残す．
追加P81の推論器滞在は10,405.1秒，この内の評価時間は26.82秒だった．時間は重なるため合算しない．Hの初回入力は全課題でMより12 token多い．全生成，API，cache，準備・作業・後片付けと評価の全件表を公開記録へ保存する．
